%% kapa: exam=1 round=20 year=2009 period=2 form=A subject=accounting subject_ko=회계학 seq=40 number=40
%% source: Q-Net 감정평가사 기출문제 (cst00309) artlSeq=5024178 file=extracted/1008/5024178/2009년 제20회 감정평가사 1차시험 기출문제/2교시.hwp
%% conversion: kordoc 4.24.51 -> tools/gen1.py (auto); review=proofread-sonnet; flags=none
\begin{question}{40}
(주)한국은 표준원가계산을 월별로 적용하고 있다. 제품 1개를 생산하기 위한 표준직접재료원가는 ￦600(4kg×\allowbreak{}@￦150)이다. 당월초에 원재료 2,400kg을 매입하였으며, 원재료 가격차이는 구매시점에서 파악한다. 당월에 생산한 제품은 500개이며, 원재료 수량차이는 ￦30,000(불리)이다. 원재료 월초재고가 200kg일 때 원재료 월말재고는 몇 kg인가?
\begin{choices}
  \item 100kg
  \item 200kg
  \item 300kg
  \item 400kg
  \item 500kg
\end{choices}
\end{question}
\answer{4}
